
%\documentclass[10pt,conference,compsocconf]{IEEEtran}
\documentclass[10pt]{article}
\usepackage{times}

\usepackage{caption}
\captionsetup{font=footnotesize,justification=centering,labelsep=period}


\usepackage{cite}


% *** GRAPHICS RELATED PACKAGES ***
%
%\ifCLASSINFOpdf
  % \usepackage[pdftex]{graphicx}
  % declare the path(s) where your graphic files are
  % \graphicspath{{../pdf/}{../jpeg/}}
  % and their extensions so you won't have to specify these with
  % every instance of \includegraphics
  % \DeclareGraphicsExtensions{.pdf,.jpeg,.png}
%\else
  % or other class option (dvipsone, dvipdf, if not using dvips). graphicx
  % will default to the driver specified in the system graphics.cfg if no
  % driver is specified.
  % \usepackage[dvips]{graphicx}
  % declare the path(s) where your graphic files are
  % \graphicspath{{../eps/}}
  % and their extensions so you won't have to specify these with
  % every instance of \includegraphics
  % \DeclareGraphicsExtensions{.eps}
%\fi



% *** MATH PACKAGES ***
%
\usepackage[cmex10]{amsmath}
\usepackage{eqnarray}


% *** SPECIALIZED LIST PACKAGES ***
%
\usepackage{algorithmic}
\usepackage[ruled, linesnumbered]{algorithm2e}
%\usepackage[justification=centering]{caption} %centrer les titres
\usepackage{amssymb}
\let\proof\relax
\let\endproof\relax



% *** ALIGNMENT PACKAGES ***
%
\usepackage{array}


\usepackage{mdwmath}
\usepackage{mdwtab}



\usepackage{eqparbox}
% Also of notable interest is Scott Pakin's eqparbox package for creating
% (automatically sized) equal width boxes - aka "natural width parboxes".
% Available at:
% http://www.ctan.org/tex-archive/macros/latex/contrib/eqparbox/




\usepackage{enumerate}

\usepackage{rotating}
\usepackage{subcaption}

\usepackage{tabularx}
\usepackage{booktabs}
\usepackage{multirow}
\providecommand{\keywords}[1]{\textbf{\textit{Index terms---}} #1}
% correct bad hyphenation here
%\hyphenation{op-tical net-works semi-conduc-tor}

%\parskip 6pt plus 2pt minus 1pt
%\parskip 3pt plus 2pt minus 1pt

\pagestyle{empty}

\begin{document}
%\pagenumbering{gobble}
%
% paper title
% can use linebreaks \\ within to get better formatting as desired
%\title{Causality based feature selection approach for multivariate time series forecasting}
\title{\textbf{\Large Chapter: Feature Selection and Dimension Reduction for Large Multivariate Time Series Forecasting}}

  
\maketitle

\begin{abstract}
Analysis, modeling and exploration of time series has recently experienced a noticeable interest gain. For a good reason, as knowledge emerging from such analysis may significantly improve the decision-making process in industry or science. For that, time series analysis is often an important part of modern Business Intelligence Systems. The hunger for information causes huge amounts of various time series to be gathered, stored and processed.

This raises some problems for traditional processing models and pushes the research in to the new directions. Utilizing forecasting models that handle large number of predictors is an increasingly popular approach now. In many cases, such models can improve results compared to univariate forecasting models. But this raises some challenges, especially in terms of Big Data.

One of the challenges is proper selection or shrinkage of predictor variables used in the forecasting model. In high dimensional time series datasets, this is a must. First, not all observed variables will have direct impact on the predicted variable, and adding unrelated variables may simply have the opposite effect on the forecasting. Secondly, there are empirical evidences (at least for the macroeconomic data) that there is little to no improvement when using too many variables.

In this chapter, we would like to discuss various aspects of selection and shrinkage of predictors in the context of massive time series datasets. While similar research exists in the literature, it is usually targeted for small and medium datasets, but to our knowledge, there is a gap in the context of Big Data.

This chapter is a high-level survey of future selection and dimension reduction tasks for time series forecasting, with a special emphasis on Big Data challenges.

\end{abstract}

\keywords{Large Time Series Forecasting, Dimension reduction. }

\section{Introduction}
Nowadays, huge amounts of numerical time series are generated by industry and research systems. For various applications, like biology, medicine, finance, industry, and many others, high dimensional time series are more and more common. Modern time series analysis systems are expected to process and store millions of such high dimensional data points per minute, 24 hours a day, seven days a week, generating terabytes of logs. As a result, the number of time series generated grows very quickly and falls under the category of Big Data (i.e. problems in which the size, dimensionality of data is a problem itself). This raises some new and interesting challenges.

Time series analysis and data mining concentrates on describing the occurring patterns and the evolution of data over the time. Among many applications, time series forecasting takes special place as it is crucial for decision making. To no surprise, it is often an important part of Business Intelligence systems. The history of forecasting models starts in the 1830s, with first univariate models. Based on the Auto-Regressive principle, data history is used in order to make forecasts. Auto-Regressive model (AR), Auto-Regressive Moving Average (ARMA), and the Auto-Regressive Integrated Moving Average (ARIMA) models \cite{box1976time}, show univariate models evolution. 

The univariate forecasting approach has some obvious drawbacks, as it skips potentially-exploitable information from other time series. To exploit this extra potential, multivariate forecasting models where developed, such as the extended version of the AR model; the Vector Auto-Regressive (VAR), and the co-integrated VAR model; the Vector Error Correction model (VECM) \cite{johansen1991estimation}. These models are still in the use nowadays, alone or combined with other new techniques like artificial neural networks dynamic modeling.

The multivariate forecasting models assume that the value of one variable often depends on past values of itself and other related variables \cite{johansen1991estimation}. In high dimensional data, this rises two main problems: (i) how to select dimensions for predicted variable, (ii) whether the number of used variables impacts predictions accuracy. Both problems are intensively studied by researchers, yielding new findings and approaches. A recent works  \cite{stock2012generalized}, \cite{jiang2017macroeconomic}, supported with empirical evidence, suggests that:
\begin{itemize}
  \item there is no selection or shrinkage method superior to others, and thus, methods should be chosen depending on data,
  \item there is a constraint on number of used variables that improves forecasting quality for existing multivariate models.
\end{itemize}


\section{Related works}
We present here some of the numerous approaches that handle the problem of forecasting time series with large number variables.

In \cite{stock2002forecasting}, Principal Components Analysis is used to forecast one time series using a large number of predictor series. A two-step forecasting model is adopted which applies principal component to reduce the number of predictors, and then a forecasting model is used.

In \cite{stock2012generalized}, various shrinkage methods are compared to the Dynamic Factor Model (DFM) \cite{geweke1977dynamic}, which is based on principal component method. The evaluation is done using a US macroeconomic time series datasets with large number of predictors.
In  \cite{jiang2017macroeconomic}, both the Australian and US data sets where used, and different shrinkage and dimension reduction methods are compared to the DFM. Authors highlighted that methods that performs well in  \cite{stock2012generalized} on US dataset does not necessarily work well with the Australian dataset. This supports the intuitive hypothesis that performance of existing models will highly depend on data, and that there is no superior method. Also in  \cite{jiang2017macroeconomic}, authors confirmed findings from  \cite{stock2012generalized}, that there is little to no improvement in forecast accuracy when the number of predictors exceeds 20-40 variables.

The Granger causality \cite{granger1980testing} has been used for feature selection, in  \cite{sun2015using}, by selecting variables that cause the target and the target does not cause them. The authors show good results compared to selected dimension reduction methods on a regression and classification models.

Dimension reduction was also used for forecasting daily stock market \cite{zhong2017forecasting}. Due to the large number of features, the authors propose a data mining forecasting process, in which they are testing three versions of principal component analysis technics as a first step, and then use Artificial Neural Network (ANN) for forecasting step.

An approach based on clustering the graph of Granger causalities between time series is proposed in  \cite{hmamouche:hal-01467523}. Evaluation is done against the ARIMA model, the Vector Error Correction Model (VECM) with all the predictors, and with some existing dimension reduction and feature selection methods. Proposed models performs well on used dataset, and outperforms other methods.

\section{Chapter proposition}

This chapter is intended as a high-level survey of feature selection and reduction tasks for time series forecasting, with a special emphasis on Big Data challenges. We plan to discuss algorithms based on DFM or dimension reduction techniques, Granger causality and Transfer entropy from the following viewpoints:
\begin{itemize}
\item scalability
\item impact of sampling and sequencing techniques
\item heterogeneous computing.
\end{itemize}

%----------------------------------------------------------------------------------------
%	BIBLIOGRAPHY
%----------------------------------------------------------------------------------------
\bibliographystyle{alpha}
\bibliography{../biblio/a,../biblio/b,../biblio/c,../biblio/d,../biblio/e,../biblio/f,../biblio/g,../biblio/h,../biblio/i,../biblio/j,../biblio/k,../biblio/l,../biblio/m,../biblio/n,../biblio/o,../biblio/p,../biblio/q,../biblio/r,../biblio/s,../biblio/t,../biblio/u,../biblio/v,../biblio/w,../biblio/x,../biblio/y,../biblio/z}

%\section*{Résumé}
%Le domaine de recherche lié au séries temporelles a connu un grand succès ces dernières décennies, et ce plus particulièrement concernant le domaine économique. De nos jours, dans un contexte de données massives (\textit{e.g.} Big Data), 
%des problèmes apparaissent avec le nombre grandissant de séries temporelles stockées car les modèles de prédiction ne peuvent utiliser toutes ces séries. En conséquence, une étape de sélection est nécessaire afin de réduire le nombre 
%de variables endogènes qui peuvent aider à la prédiction de chaque série temporelle.  
%
%
%L'approche standard consiste à appliquer un modèle de prédiction en utilisant toutes les autres séries à disposition, ou celle ayant la plus grande causalité vis à vis de la série à prédire. 
%Même si cette approche  semble évidente, elle ne prend pas en compte les relations cachées entre les variables.  
%
%
%Dans cet article, nous proposons la méthode GSM (Granger Selection Method) pour sélectionner $n$ attributs à partir d'un ensemble de séries temporelles. Notre méthode est basée sur
%($i$) le graphe de causalité de Granger et ($ii$) une adaptation de l'algorithme $k$-medoïde, une autre méthode de partitionnement, afin de construire des ensembles de séries temporelles. 
%Finalement, des expérimentations menées sur des jeux de données réels confirment notre approche.


\end{document}
